\section{Jai Hanuman}
\label{sec:g-jai-hanuman}
\problemlink{https://www.codechef.com/problems/ACM14KG4}{CodeChef ACM14KG4}
\source{ICPC 2014--15 Kharagpur Online Re-contest (PYQ, ACM14KG4)}{Medium--Hard}

\problemstatement{$N$ people and $5N$ crackers numbered $1..5N$; person $i$
starts with crackers $3i-2,3i-1,3i$. In order $1,2,\dots,N$ each person may
make up to $3$ swaps with anyone (swaps cannot be refused, Hanuman's spare
crackers included). Everyone wants three \emph{consecutive} crackers; among
those, the least one should be $\ge P_i$ and as small as possible; if that is
impossible, the least one should be as small as possible. Everybody plays
optimally and knows everything. Print each person's least cracker.}

\begin{patternbox}
Sequential moves where the \emph{last} mover cannot be disturbed
$\Rightarrow$ \textbf{backward induction}: fix the last player's outcome,
then the second last's among what is left, and so on.
\end{patternbox}

\subsection*{Approach and intuition}

Three swaps are enough to collect any three crackers, so person $N$, who moves
last, always ends with exactly the triple they want. Person $N-1$ knows this,
so anything they take from $N$'s target set will be taken away again; their
best achievable triple is the best one \emph{avoiding} $N$'s triple. Inductively,
person $i$ gets their preferred triple among crackers not wanted by
$i+1,\dots,N$.

Implementation: a start $s$ is \emph{available} while $s,s+1,s+2$ are all
free. Keep all available starts in a \texttt{set}. For person $i$ (from $N$
down to $1$) take \texttt{lower\_bound}$(P_i)$, or the smallest start if none
exists. Taking $[s,s+2]$ kills starts $s-2,\dots,s+2$. At most $5$ starts die
per person and there are $5N-2$ initially, so a start always remains.

\subsection*{Algorithm}
\begin{algorithmic}[1]
\State $S\gets\{1,\dots,5N-2\}$
\For{$i=N$ \textbf{down to} $1$}
  \State $s\gets$ smallest element of $S$ that is $\ge P_i$, else $\min S$
  \State $ans_i\gets s$; remove $s-2..s+2$ from $S$
\EndFor
\end{algorithmic}

\dryrun{$N=2$, $P=(4,9)$, crackers $1..10$. Person~2: starts $\ge9$ need
$s\le8$, none, so $s=1$ (crackers $1,2,3$). Person~1: smallest free start $\ge4$
is $4$. Output \texttt{4 1}. With $P=(1,5)$: person~2 takes $5$, person~1 takes
$1$: \texttt{1 5}. \checkmark}

\subsection*{C++17 implementation}
\cppfile{greedy/26_jai_hanuman.cpp}

\begin{complexitybox}
\textbf{Time:} $O(N\log N)$ per test. \textbf{Space:} $O(N)$.
\end{complexitybox}

\begin{pitfallbox}
\begin{itemize}
  \item Processing people from $1$ to $N$ (forward) is the natural but
        wrong reading --- earlier players do not get to keep contested
        crackers.
  \item Consecutive means \emph{any} three consecutive numbers, not aligned
        blocks $\{3i-2,3i-1,3i\}$.
\end{itemize}
\end{pitfallbox}
